% Generated draft: requires mathematical and source review.
\documentclass[11pt]{article}
\usepackage[margin=25mm]{geometry}
\usepackage{fontspec}
\setmainfont{Latin Modern Roman}
\usepackage{amsmath,amssymb,mathtools,unicode-math}
\setmathfont{Latin Modern Math}
\usepackage{enumitem,needspace,xurl,hyperref}
\setlength{\parindent}{0pt}
\setlength{\parskip}{5pt}
\setlength{\emergencystretch}{3em}
\newcommand{\N}{\mathbb{N}}
\newcommand{\Z}{\mathbb{Z}}
\newcommand{\Q}{\mathbb{Q}}
\newcommand{\R}{\mathbb{R}}
\newcommand{\C}{\mathbb{C}}
\newcommand{\F}{\mathbb{F}}
\newcommand{\A}{\mathbb{A}}
\newcommand{\PP}{\mathbb{P}}
\newcommand{\E}{\mathbb{E}}
\newcommand{\eps}{\varepsilon}
\newcommand{\dd}{\,\mathrm{d}}
\begin{document}
\section*{UnsolvedMath ID: 20000376}
\textbf{Source problem code:} catalogue entry
\subsection{Definitions and mathematical statement}
Let \(k\) be a number field, \(\bar k\) an algebraic closure, and \(X/k\) a K3 surface, understood to be smooth, projective and geometrically integral, with trivial canonical bundle and \(H^1(X,\mathcal O_X)=0\). Put
\[
 \operatorname{Br}(X)=H^2_{\mathrm{\acute et}}(X,\mathbb G_m)_{\mathrm{tors}},\quad
 \operatorname{Br}_0(X)=\operatorname{im}(\operatorname{Br}(k)\to\operatorname{Br}(X)),\quad
 \operatorname{Br}_1(X)=\ker(\operatorname{Br}(X)\to\operatorname{Br}(X_{\bar k})).
\]
For an abelian group \(B\), write \(B[n]=\{b:nb=0\}\). For \(\alpha\in\operatorname{Br}(X)\), its Brauer--Manin pairing with an adelic point is \(\sum_v\operatorname{inv}_v\alpha(P_v)\). An obstruction to the Hasse principle means that \(X(\A_k)\ne\varnothing\), but the points orthogonal to the specified class or subgroup form the empty set; it then follows that \(X(k)=\varnothing\).

The 2014 question attributed by V\'arilly-Alvarado to Skorobogatov asks whether there exist a K3 surface \(X/\Q\) and an odd-order class in \(\operatorname{Br}(X)\) giving such an obstruction, and whether this can already occur in \(\operatorname{Br}_1(X)\).

For the quartic specialization \(X\subset\PP^3_{\Q}\), the source fixes an odd positive integer \(u\) and \(\alpha\in\operatorname{Br}(X)[u]\). It states that, for every place \(v\), there exists a zero-cycle \(Z_v\) on \(X_{\Q_v}\) of degree one orthogonal to \(\alpha\). Here a zero-cycle is a finite integral linear combination of closed points, its degree is the corresponding sum of residue-field degrees, and orthogonality locally means that the invariant of the sum of the corestricted evaluations of \(\alpha\) is zero. The source then invokes an unspecified conjecture of Colliot-Th\'el\`ene to obtain a global degree-one zero-cycle and asks whether a rational point follows. The precise conjecture and the passage from orthogonality for a single class to the hypotheses of that conjecture are not supplied.

The explicit quartic question is whether every smooth quartic K3 surface \(X/\Q\) satisfying
\[
 \bigl(\operatorname{Br}(X)/\operatorname{Br}_0(X)\bigr)[2]=0
\]
satisfies \(X(\A_{\Q})\ne\varnothing\Rightarrow X(\Q)\ne\varnothing\). This retains the source's quotient \(\operatorname{Br}(X)/\operatorname{Br}(\Q)\), interpreted using the image of the constants. Tschinkel also asks about weak approximation, meaning density of \(X(\Q)\) in finite products \(\prod_{v\in S}X(\Q_v)\) for finite sets of places. Skorobogatov's response that this will probably fail is a prediction, not a theorem in the record. Hassett asks for the corresponding degree-six situation \(X\subset\PP^4\); the text does not restate the hypotheses or specify precisely which preceding implication is intended in that variant.

The supplied later research report describes affirmative algebraic and transcendental three-torsion examples on degree-two K3 surfaces, while retaining the quartic and degree-six questions as unresolved in the literature it checked. These reports do not identify degree-two examples with the quartic or degree-six cases. Its separately stated candidate lemma is: for any polarized K3 surface \((X,h)\) over a number field and any prime \(p\),
\[
 (\operatorname{Br}_1(X)/\operatorname{Br}_0(X))[p]\ne0
 \quad\Longrightarrow\quad
 \rho(X_{\bar k})=\operatorname{rank}\operatorname{Pic}(X_{\bar k})\ge p.
\]
The report additionally claims a simultaneous prime-to-embedding-degree construction of adelic degree-one zero-cycles orthogonal to the full Brauer group, with a conditional global zero-cycle conclusion for quartics and a remaining three-primary issue in degree six. Its full hypotheses and proof are absent; this is retained as a reported contribution requiring review, rather than an unconditional replacement for the original questions. A zero-cycle of degree one need not itself be a rational point.
\subsection{Short English statement}
Can odd-order Brauer classes, even algebraic ones, obstruct rational points on a K3 surface over \(\Q\)? The source also asks whether quartic K3 surfaces with no two-torsion in the Brauer quotient satisfy the Hasse principle, and raises weak approximation and degree-six variants. Later supplied reports claim three-torsion examples in degree two and a Picard-rank bound, but do not settle the stated quartic question or justify replacing degree-one zero-cycles by rational points.
\subsection{Sources}
\begin{itemize}
\item Local UnsolvedMath registry, permanent ID \(20000376\), source code \texttt{catalogue entry}; AIM Problem Lists, algebraic number theory, canonical \texttt{aim-algebraic-number-theory-notes.json}, notes[27]. Attribution: Ulam AI / UnsolvedMath Contributors, CC BY 4.0.
\item AIM workshop, \emph{Rational and integral points on higher-dimensional varieties}, Problem 8, May 28, 2014: \url{https://aimath.org/pastworkshops/ratlhigherdimvarproblems.pdf}. The dialogue names V\'arilly-Alvarado, Skorobogatov, Tschinkel and Hassett, and refers to a conjecture of Colliot-Th\'el\`ene without a bibliographic citation.
\item The supplied verification labels the statement exact and reports checking the four-page PDF and font encoding; in particular \(u\) in \(\operatorname{Br}(X)[u]\) is reported to be the literal letter \(u\), with \(u\) odd. This draft relies on that supplied account rather than a new source inspection.
\item Registry record dated August 14, 2026: \texttt{partially\_solved}, \texttt{PARTIAL-PROGRESS}; research attempt 1 is labelled partial result, rigor proved, novelty confidence low. Its degree-two examples, Picard-rank lemma and zero-cycle construction are machine-generated research aids, not peer reviewed, and require independent verification before citation. No primary citations for these later claims are supplied.
\end{itemize}
\end{document}
